\chapter{Textes grecs et bilingues}\label{ch:grec}

Sept textes de l’ensemble C comportent une partie grecque, seule ou accompagnée de latin : \nnums{8}, 21, 110, 168, 348, 360 et 447 (tableau~\ref{tab:grecs}). Aucune de ces parties grecques ne figurait dans l’export de 2020 de l’EDCS, qui omettait les caractères grecs ; elles ont été récupérées dans l’état 2022 de la base et collationnées sur les fiches de l’EDB. Le groupe est trop petit pour une étude quantitative. Il apporte quatre séries d’observations.

% généré par scripts/04_generated.py
\begin{longtable}{@{}r r p{0.78\linewidth}@{}}\caption{Textes grecs et bilingues (texte EDCS, état 2022).}\label{tab:grecs}\\\toprule N° & Année & Texte\\\midrule\endfirsthead\bottomrule\endfoot
8 & 343 & {\small κατάθεσις / Ὀκταβίλλης / πρὸ θʹ καλ(ανδῶν) σε(π)τ(εμβρίων) / (ὑπατείᾳ) Ῥωμουλι} \\
21 & 347 & {\small ] quae fecit cum m[arito ann(os) - - -] / [quae vixit a]nn(os) XXXV m(enses) II d(ies) X d[ep(osita) - - -] / [- - - Ἡ]λιοδώρα ἐτῶν τρ[ιάκοντα πέντε] / [ὑπατείᾳ Ῥουφίνου καὶ Εὐ]σεβίου μηνὶ μαΐ[ου τό - - -]} \\
110 & 367 & {\small ἐν εἰρήνῃ ἀνεπάη Βεναινᾶτος ἀπ(ὸ) κο(μήτων) / ζήσας ἔτη οεʹ ὃς κῖτε ἐν τῷδε τόπῳ ἐξελτὼν / ἐκ τοῦ βίου τῇ πρ(ὸ) δʹ ἰδ(ῶ)ν μαρ(τίων) ὑπ(ατείᾳ) Λουππικί(νου) / καὶ Ἰοβίνου Θεοφίλα ἐποίησα συμβίῳ / Benenato Theofila fecit uxor eius} \\
168 & 373 & {\small [Locus - - - u]bi posita es[t - - -] / [Valente Aug(usto)] IIII et Valentinia[no Aug(usto) IIII conss(ulibus)] / [- - -] in pace qu(a)e vixit annis XXV[- - -] / [τόπος - - -]νου ζῶντος ἐν ᾧ κατάκιτε / πεδ[ὶν Λέων(?) κ]αὶ εὐγενεστάτη Κ[ω]νσταντινου / πολιτί[σσα Βι]ταλία ἐτῶ[ν εἴκο]σι ἑπτά} \\
348 & 392 & {\small ἐνθάδε κατάκι\{τι\}τε Ἀλέξα(ν)δρο(ς) ζήσας / ἔτη εἴκοσι τελευτᾷ μηνὶ νοενβρί / ου υ\{ε\}ἱὸς Ἀ(μ)βρωσίου (χ)ωρίῳ Μικρᾶς Κ / ώμης ζέσας μ(ε)τὰ τῶν ἁ(γ)ίων} \\
360 & 393 & {\small ] / [- - - ἔ]τη ιʹ Θεο[δοσίῳ τό - - -] / [} \\
447 & 401 & {\small ὁρκίζομε δὲ ὑμᾶς [τὸν θ(εὸ)ν τὸν πα] / ντοκράτορα εἴ τις το[λμήσει σκυ] / λῖσε τὰ ὀστέα αὐτοῦ [ἢ ἄλλον τινὰ] / [θῖνε ἔσ]τε ἔνοχος τῇ [κρίσει - - -] / [- - -]αν ἡμαίραν τῆς [ἀναστάσεως(?)]} \\
\end{longtable}


\section{La datation romaine transposée en grec}

La formule consulaire est rendue par \emph{ὑπατείᾳ} suivi des noms au génitif (\nnums{8}, 21, 110 ; \emph{ὑπατίᾳ} dans le texte restitué par l’EDB au \nnum{447}). Le calendrier romain est transposé de la même manière : \emph{πρὸ θʹ καλ(ανδῶν) σε(π)τ(εμβρίων)} au \nnum{8}, \emph{τῇ πρ(ὸ) δʹ ἰδ(ῶ)ν μαρ(τίων)} au \nnum{110}. Le \nnum{348} donne au contraire un mois sans quantième (\emph{μηνὶ νοενβρίου}), et le \nnum{21} un mois suivi d’un quantième perdu.

Au \nnum{8}, l’épitaphe d’Octavilla, le mot \emph{ὑπατείᾳ} n’est pas gravé : l’EDB le place entre chevrons, signe d’une omission du lapicide suppléée par l’éditeur, et l’EDCS entre parenthèses. Le texte gravé porte seulement \emph{Ῥωμουλι}, nom au génitif du consul de 343. La notice le signale ; l’année n’en est pas affectée.

\section{Les graphies}

Les graphies témoignent de la prononciation contemporaine. La monophtongaison de \emph{αι} en \emph{ε} et celle de \emph{ει} en \emph{ι} sont constantes : \emph{κῖτε} pour \emph{κεῖται} (\nnum{110}), \emph{κατάκιτε} pour \emph{κατάκειται} (\nnums{168} et 348), \emph{ὁρκίζομε} pour \emph{ὁρκίζομαι} et \emph{σκυλῖσε} pour \emph{σκυλῖσαι} (\nnum{447}) ; à l’inverse, \emph{ἡμαίραν} pour \emph{ἡμέραν} (\nnum{447}) est une graphie inverse. Le \nnum{110} écrit \emph{ἐξελτών} pour \emph{ἐξελθών}, et le nom latin \emph{Benenatus} y devient \emph{Βεναινᾶτος}. Au \nnum{348}, \emph{νοενβρίου} transcrit le latin \emph{novembris}.

\section{Identités et milieux}

Trois textes renseignent sur l’origine ou le statut des défunts. Au \nnum{110}, Benenatus est dit \emph{ἀπ(ὸ) κο(μήτων)}, « ancien \emph{comes} » ; son épouse Theophila fait graver une épitaphe bilingue, la partie latine se réduisant à la dédicace (\emph{Benenato Theofila fecit uxor eius}). Au \nnum{168}, l’emplacement acquis du vivant d’un homme dont le nom est perdu (\emph{[…]νου ζῶντος}) accueille un enfant et une femme qualifiée d’\emph{εὐγενεστάτη Κ[ω]νσταντινουπολίτ[ισσα]} ; le nom de l’enfant (\emph{Λέων ?}) et une partie de l’ethnique sont restitués. Au \nnum{348}, Alexandros, fils d’Ambrosios, est originaire d’un village (\emph{χωρίῳ}) nommé Mikra Kômè, que le texte ne situe pas.

Selon l’EDB, la première partie du \nnum{447}, absente de l’EDCS, nomme un défunt dont seul le début du nom est conservé (\emph{Εὐ[…]}) et le qualifie de \emph{[Θεο]δοσειοπολίτ[ης]}, originaire de Théodosiopolis. Le nom de la ville est en partie restitué, et plusieurs villes ont porté ce nom à la fin du \textsc{iv}\textsuperscript{e} siècle : l’identification ne peut être précisée sur la seule foi du texte.

\section{Deux textes de portée religieuse}

\emph{Le \nnum{348}.} Le texte se termine par \emph{ζέσας μ(ε)τὰ τῶν ἁ(γ)ίων}. La forme transmise \emph{ζέσας} n’est pas grammaticalement satisfaisante dans ce contexte, puisque le participe \emph{ζήσας} figure déjà plus haut pour la durée de vie. Elle peut être comprise comme une graphie phonétique de l’optatif de souhait \emph{ζήσαις}, « puisses-tu vivre avec les saints ». Cette lecture rattacherait le texte au type, bien attesté dans l’épigraphie chrétienne grecque, de l’acclamation au défunt ; elle reste une interprétation, qu’il faudrait confronter à l’édition et à la photographie. L’année 392 du \nnum{348} provient en outre de la seule base : le texte ne contient aucune formule consulaire (chapitre~\ref{ch:fondement}).

\emph{Le \nnum{447}.} Le texte contient une formule d’adjuration contre la violation de la tombe : \emph{ὁρκίζομε δὲ ὑμᾶς […]ντοκράτορα εἴ τις το[λμήσει σκυ]λῖσε τὰ ὀστέα αὐτοῦ}, « je vous adjure par [Dieu tout-puis]sant : si quelqu’un ose déranger ses os… ». Le nom de Dieu est entièrement restitué ; seule la fin de \emph{[παντο]κράτορα} est lue. La sanction est restituée différemment par les deux bases. L’EDCS donne \emph{ἔνοχος τῇ [κρίσει …] […]αν ἡμαίραν τῆς [ἀναστάσεως ?]}, « passible du jugement […] au jour de la résurrection » ; l’EDB restitue \emph{ἔνοχος τῇ [κολάσει τὴν φοβερ]ὰν ἡμαίραν τῆς [κρίσεως ἔστω]}, « qu’il soit passible du châtiment au jour redoutable du jugement ». Dans les deux cas, les mots \emph{κρίσις} et \emph{ἀνάστασις} sont restitués. Le texte transmis atteste une menace eschatologique (\emph{ἔνοχος}, \emph{ἡμαίραν τῆς}) ; il n’atteste ni le mot « résurrection » ni le mot « jugement ». La version précédente du manuscrit traduisait la restitution de l’EDCS sans la signaler comme telle ; la notice et le chapitre~\ref{ch:mort_esperance} en tiennent désormais compte.

Ces deux textes sont les seuls du groupe grec à porter un énoncé religieux au-delà du formulaire. Avec l’\emph{ἐν εἰρήνῃ ἀνεπάη} du \nnum{110}, équivalent du \emph{requiescit in pace} latin, ils montrent que la communauté hellénophone de Rome emploie, dans les cimetières de la ville, un formulaire parallèle au formulaire latin.
